% This is the ADASS_template.tex LaTeX file, 24 Oct 2025.
% It is based on the ASP general author template file, but modified to reflect the specific
% requirements of the ADASS proceedings.

% This is the aspauthor.tex LaTeX file
% Copyright 2021, Astronomical Society of the Pacific Conference Series
% Revision:  29 November 2021

% To compile, at the command line positioned at this folder, type:
% pdflatex ADASS_template
% pdflatex ADASS_template
% This will create a file called ADASS_template.pdf

\documentclass[11pt,twoside]{article}

% Do NOT use ANY packages other than asp2023. 
\usepackage{asp2023}

\aspSuppressVolSlug
\resetcounters

% References must all use BibTeX entries in a .bibfile.
% References must be cited in the text using \citet{} or \citep{}.
% Do not use \cite{}.
% See ManuscriptInstructions.pdf for more details
\bibliographystyle{asp2021}

% The ``markboth'' line sets up the running heads for the paper.
% 1 author: "Surname"
% 2 authors: "Surname1 and Surname2"
% 3 authors: "Surname1, Surname2, and Surname3"
% >3 authors: "Surname1 et al."
% Replace ``Short Title'' with the actual paper title, shortened if necessary.
% Use mixed case type for the shortened title
% Ensure shortened title does not cause an overfull hbox LaTeX error
% See ASPmanual2010.pdf 2.1.4  and ManuscriptInstructions.pdf for more details
\markboth{Danehkar et al.}{AstroAKD for Astronomical Archive Search}

\begin{document}

\title{AstroAKD: An MCP-Based Agent for Natural-Language Search of Astronomical Archives}

% Note the position of the comma between the author name and the 
% affiliation number.
% Authors surnames should come after first names or initials, eg John Smith, or J. Smith.
% Author names should be separated by commas.
% The final author should be preceded by "and".
% Affiliations should not be repeated across multiple \affil commands. If several
% authors share an affiliation this should be in a single \affil which can then
% be referenced for several author names. If only one affiliation, no footnotes are needed.
% See ManuscriptInstructions.pdf and ASP's manual2010.pdf 3.1.4 for more details
\author{Ashkbiz~Danehkar,$^1$ Gaurab~Panthee,$^2$ Pushwitha Krishnappa,$^2$ Nishan~Pantha,$^2$ Nidhi Jha,$^2$ Muthukumaran Ramasubramanian,$^2$ and Rahul Ramachandran$^3$}
\affil{$^1$Universities Space Research Association, Huntsville, AL, 35805, USA \newline %;%\email{corresponding.author@example.org}}
$^2$University of Alabama in Huntsville, Huntsville, AL, 35899, USA \newline
$^3$NASA Marshall Space Flight Center, Huntsville, AL 35812, USA
}

% This section is for ADS Processing.  There must be one line per author. paperauthor has 9 arguments.
\paperauthor{Ashkbiz Danehkar}{adanehkar@usra.edu}{0000-0003-4552-5997}{Universities Space Research Association}{Science and Technology Institute}{Huntsville}{AL}{35805}{USA}
\paperauthor{Gaurab Panthee}{gp0065@uah.edu}{0009-0009-7292-8960}{University of Alabama in Huntsville}{Computer Science}{Huntsville}{AL}{35899}{USA}
\paperauthor{Pushwitha Krishnappa}{pk0055@uah.edu}{0009-0009-2482-1149}{University of Alabama in Huntsville}{Computer Science}{Huntsville}{AL}{35899}{USA}
\paperauthor{Nishan Pantha}{np0069@uah.edu}{0009-0003-6948-1463}{University of Alabama in Huntsville}{Earth System Science Center}{Huntsville}{AL}{35899}{USA}

% There should be one \aindex line (commented out) for each author. These are used to
% build up the author index for the Proceedings. The surname must come first, followed by
% initials. Note the use of ~ before each initial to control spacing.
% The \author entries (see above) have surname last. These \aindex entries have
% surname first.
% The Aindex.py command willl create them for you after you have constructed the \author
% The first entry should be the first author, for bold-facing the author index page-reference

%\aindex{FistAuthor1,~S.~A.}
%\aindex{Author2,~S.~B.}
%\aindex{Author3,~S.}


\begin{abstract}
Astronomical data search requires object identification, archive selection, construction of service-specific queries, and interpretation of heterogeneous metadata. We present AstroAKD, a natural-language search agent developed within NASA IMPACT's Accelerated Knowledge Discovery (AKD) ecosystem. An Astroquery-based Model Context Protocol (MCP) server exposes astronomical archive functions to an LLM agent while retaining executable queries and retrieved evidence. We also present AstroAKD-Bench, which generates reference-aware synthetic questions, collects independent MCP reference responses, and evaluates three GPT-5.2 configurations: MCP-only, MCP with AstroAKD instructions, and web-only search. A set of 100 synthetic questions comprised 60 archive-driven, 20 coordinate-driven, 12 multi-archive, and 8 literature-driven requests. Complete reference evidence enabled automated scoring for 92 questions; eight literature cases were excluded. Strict agreement with the archive-derived references was 94.57\%, 95.65\%, and 14.13\% for the three configurations, respectively. We separately report execution completion, latency, token usage, and trace provenance. These results quantify reference agreement, rather than independently verified scientific accuracy.
AstroAKD demonstrates how domain-specific agents equipped with MCP tools can support rapid exploration of data archives.
%Astronomical data search requires the identification of objects, selection of appropriate data archives, construction of archive-specific queries, and interpretation of heterogeneous metadata. We present AstroAKD, a natural-language data-search agent developed within NASA IMPACT's Accelerated Knowledge Discovery (AKD) ecosystem. An Astroquery-based Model Context Protocol (MCP) server exposes supported astronomical archive functions to an LLM agent while preserving executable queries and returned evidence. We also present AstroAKD-Bench, which generates synthetic queries, collects independent MCP reference responses, and compares three agent configurations: (A) MCP only, (B) MCP plus AstroAKD instructions, and (C) web search only. For 100 generated questions, 92 were evaluable across 80 single-call and 12 multi-archive tasks, while 8 literature tasks were excluded. For a set of generated queries, agreements with these reference snapshots were 94.57\% (A), 95.65\% (B), and 18.84\% (C). We report operational and reference-agreement measures separately and retain queries and traces for reproducibility. AstroAKD demonstrates how domain-specific agents equipped with MCP tools can support rapid exploration of data archives.
\end{abstract}

% These lines show examples of subject index entries. At this stage these have to commented
% out, and need to be on separate lines. Eventually, they will be automatically uncommented
% and used to generate entries in the Subject Index at the end of the Proceedings volume.
% Don't leave these in! - replace them with ones relevant to your paper.

% These lines show examples of ASCL index entries. At this stage these have to commented
% out, and need to be on separate lines. Eventually, they will be automatically uncommented
% and used to generate entries in the ASCL Index at the end of the Proceedings volume.
% The ascl.py command will scan your paper on possible code names.
% Don't leave these in! - replace them with ones relevant to your paper.
%\ooindex{FOOBAR, ascl:1101.010}


\section{Introduction}
Astronomical data search spans multi-mission archives and catalog services. A research query can involve identifying an object, resolving its coordinates, selecting a data archive, constructing archive-specific parameters, filtering the returned metadata, and iteratively refining the request. The Python package {\tt Astroquery} \citep{Ginsburg2019} provides common programmatic access to many such services, but users must still select the appropriate archives and translate scientific requests into service-specific parameters.

Large language models (LLMs) are increasingly used for astronomy-related information extraction, literature review, and research \citep[e.g.,][]{Grezes2024,Iyer2024,Shao2024,Ting2025}. An LLM's internal knowledge alone is insufficient as an authoritative source of observational facts. A useful scientific agent should instead plan executable archive requests and derive answers from externally returned evidence.

\textit{AstroAKD} is an astronomical data-search agent developed within NASA IMPACT's Accelerated Knowledge Discovery (AKD) ecosystem.\footnote{\url{https://github.com/NASA-IMPACT/akd-suite}} AKD provides a modular environment for developing and assessing scientific agents with inspectable execution traces. AstroAKD connects an LLM to external services through the Model Context Protocol \citep[MCP, e.g.,][]{Hou2026,Stratis2026}, which exposes discoverable tools with structured arguments. This separation supports systematic examination of archive routing, tool arguments, returned records, and answer formation.

\section{System Architecture}

Figure~\ref{fig:architecture} summarizes the system architecture. Given a natural-language request, the agent interprets object identity and observational constraints, discovers supported capabilities, and selects functions exposed by the MCP server.\footnote{\url{https://github.com/NASA-IMPACT/astroquery-mcp}} The server adapts structured parameters to {\tt Astroquery} calls, retrieves records from external services, and returns a machine-readable representation. The agent then constructs its response from that evidence rather than from unverified internal knowledge.

\articlefigure[width=.96\textwidth]{fig1_AstroAKD.eps}{fig:architecture}{AstroAKD architecture. The agent plans the search; the MCP layer exposes {\tt Astroquery} functions and returns structured evidence and inspectable traces.}

%\subsection{Astroquery as the archive-access layer}


\noindent \textbf{Astroquery as the archive-access layer.} {\tt Astroquery} \citep{Ginsburg2019} is an Astropy-affiliated set of Python clients for astronomical archives. It wraps REST APIs, TAP services, and related protocols rather than requiring AstroAKD to reimplement archive access. The broader MCP configuration includes object-resolution and catalog services (e.g., SIMBAD, VizieR), mission archives (e.g., MAST, HEASARC, IRSA), Gaia and SDSS, ESA services, the Exoplanet Archive, and NASA ADS/SciX literature search. {\tt Astroquery} uses {\tt Astropy} \citep{AstropyCollaboration2022} for coordinates, physical quantities, and related representations.

%\subsection{Dynamic MCP interface}
\smallskip

\noindent \textbf{Dynamic MCP interface.} The MCP server avoids defining a unique agent tool for every archive function. It allows the agent to list configured modules, discover available functions, inspect signatures and parameter descriptions, and execute selected functions. Structured arguments are adapted to astronomy-specific quantities, including sky positions and angular radii. Returned tables include records, column metadata, and units where provided. Errors are retained in structured traces so agents and evaluators can distinguish service failures from missing or invalid information.


\section{Agent Workflow and Evaluation}

A multi-service request may begin with the coordinates of a named target, continue with a search of observational archives, and incorporate literature from ADS. Each step maps to explicit MCP calls that can be inspected. We assess the complete agent behavior as well as the final answer using \textit{AstroAKD-Bench}, a separate benchmark built to generate questions and obtain reference answers through the Astroquery MCP server.\footnote{\url{https://github.com/NASA-IMPACT/astroakd-bench}}

AstroAKD-Bench generated 100 synthetic reference-aware questions, including 60 archive-driven, 20 coordinate-driven, 12 multi-archive, and eight literature-driven tasks. Independent Astroquery MCP calls supplied complete reference evidence for all 80 single-call tasks and all 12 multi-archive tasks; the 8 literature-driven tasks remained excluded because verified publication-to-observation associations were not available. The same GPT-5.2 model and original execution limits were used in three configurations: (A) a general-purpose agent with Astroquery MCP tools, (B) the MCP-enabled agent augmented with AstroAKD domain instructions, and (C) web-only search without MCP. Answers were scored by comparing relevant numeric values, coordinates, structured fields, and duplicate-sensitive lists against archive-returned reference values, with 92 questions forming the common denominator. 
Reference answers were collected through separate Astroquery MCP executions, independently of the evaluated agent responses. Strict matching measures agreement with archive-derived reference snapshots rather than independently verified scientific correctness. 

The resulting strict reference-agreement rates were 87/92 (94.57\%), 88/92 (95.65\%), and 13/92 (14.13\%) for A, B, and C, respectively (Table~\ref{tab:bench}). Operational completion was 98/100, 100/100, and 100/100 in the original common-budget runs. Two failed MCP-only executions were subsequently completed in a separate higher-turn diagnostic, which is not mixed into the original operational comparison. The small one-question difference between A and B should not be interpreted as a demonstrated improvement before remaining scoring ambiguities are resolved. Because the benchmark emphasizes astronomical archive queries and uses Astroquery-derived references, the comparison is intended to evaluate archive-query capabilities rather than general-purpose web-search performance.

%A set of 100 reference-aware synthetic queries were created using AstroAKD-Bench: 60 archive-driven, 20 coordinate-driven, 8 literature-driven, and 12 multi-archive questions. Candidate single-call queries were screened for bounded, manageable, nonempty responses before question authoring.   Twenty literature or multi-archive questions require separate multi-step reference evidence; of the 80 single-call cases, 11 have answer structures unsupported by the current exact-value/coordinate matcher. Thus, 92 queries are scorable, which include question-aware scoring and multi-archive collection, while eight literature queries remain outstanding,

%Three cases used GPT-5.2 and a common baseline reasoning-turn budget: (1) a general agent with Astroquery MCP tools and no domain instruction; (2) the same MCP access augmented by AstroAKD/Collaborative Agent Reasoning Engineering \citep[CARE;][]{Ramachandran2026} guidance; and (3) web search without MCP access. The benchmark saved each run's response, status, tool trace, runtime, and token use. Agent answers were compared against i%ndependent reference snapshots using exact scalar or complete single-field list matching, coordinate comparison within $10^{-5}$ degrees, 
%question-aware archive/field-associated tuples, numeric tolerances, duplicate-sensitive list matching, and coordinate tolerance. Failed executions and abstentions count as nonmatches on scorable questions. The comparison is automatic reference-snapshot agreement. %, not independently reviewed astrophysical accuracy.




\begin{table}[!ht]
\caption{Measured AstroAKD-Bench results for synthetic queries.}
\label{tab:bench}
\smallskip
\begin{center}
{\small
\begin{tabular}{lccc}
\tableline
\noalign{\smallskip}
\textbf{Measure} & \textbf{MCP only} & \textbf{MCP+Instruction} & \textbf{Web only} \\
\noalign{\smallskip}
\tableline
\noalign{\smallskip}
Strict matches ($N/92$) & 87/92 & 88/92 & 13/92 \\
Reference agreement (\%) & 94.57 & 95.65 & 14.13 \\
%Completed runs, $N/100$ & 98/100 & 100/100 & 100/100 \\
Partially correct ($N/92$) & 4 & 4 & 2 \\
Mean runtime (s/run) & 21.20 & 21.33 & 16.19 \\
Total recorded tokens & 800,025 & 1,474,024 & 1,719,471 \\
\noalign{\smallskip}
\tableline
\end{tabular}}
\end{center}
\noindent{\footnotesize \textbf{Note.} 
92 evaluable queries include 80 single-call and 12 multi-archive; 8 literature tasks remain unscored. Partially correct answers match some reference evidence but do not satisfy complete-answer requirements and are excluded from strict matches. Runtime and token usage exclude higher-turn diagnostic retries. %are from the original common-budget runs, excluding higher-turn diagnostic retries.
}
\end{table}

%Table~\ref{tab:bench} summarizes the reference-answer agreement and operational results. The MCP-enabled configurations produced more automatically matched answers than web search on this subset, whereas domain guidance increased by one strict match among 92. 
%did not increase the strict match fraction relative to MCP-only access. 
%The web-only arm was faster on average under this setup, while its total token use was highest. %The MCP-only baseline had two execution errors (literature-driven questions) at the original eight-turn limit. Both completed in a separate diagnostic retry. %; those changed-budget runs are not substituted into Table~\ref{tab:bench}.

\section{Discussion and Conclusions}

AstroAKD treats an LLM as a planner and reasoning layer rather than an authoritative source of astronomical observations. MCP provides a standardized tool boundary and {\tt Astroquery} supplies archive-specific clients. %The AstroAKD-Bench results demonstrate the use of direct archive-tool access for agreement with automatically collected reference responses. The two comparable MCP-based agents match rates caution against equating a more elaborate prompt with higher scientific correctness without a defined evaluation rubric.
The AstroAKD-Bench evaluation shows that both MCP-enabled configurations achieve substantially higher agreement with archive-derived reference answers than the web-only baseline on the evaluated synthetic questions. The difference between the MCP-only and AstroAKD-instructed configurations is small, and further task-level analysis is needed to establish the contribution of domain guidance.
Because the benchmark uses Astroquery-derived references, the comparison primarily measures the agents' ability to retrieve and reproduce structured astronomical archive data, rather than their general question-answering capability.


The automatic comparison quantifies agreement with independently collected archive responses, not externally audited astrophysical truth. The latest scoring policy supports 92 of the 100 synthetic questions; the remaining eight are literature-driven tasks needing verified links between publications and observations. Output normalization, response completeness, archive-to-field associations, and minor identifier-format ambiguities remain evaluation sensitivities. The 12 multi-archive references were recovered using read-only, resumable MCP collection, while a separate diagnostic resolved two agent execution failures under an increased reasoning-turn budget. Future work will complete literature evidence collection, audit borderline scoring cases, and strengthen provenance and cross-archive result validation without regenerating the benchmark.


%There are limitations. Reference snapshots are archive-derived; the present strict comparator supports only 92 of the 100 questions; 80 single-call references and 12 complete two-archive reference records were available, providing 92 evaluable questions; eight literature-driven questions remained excluded. The multi-step literature questions require separate evidence, while additional answer types need appropriate structured comparison rules. Large responses, ambiguous targets, rate limits, and service unavailability remain practical challenges. Thus, the reported match rates should not be generalized to all astronomy archive tasks or interpreted as validated astrophysical accuracy. Future work will improve multi-step reference collection, structured answer scoring, bounded response handling, and error recovery while retaining query and trace provenance.

\acknowledgements This work was supported by NASA Grant 80MSFC22M004.

\bibliography{C104}  % For BibTex

% if we have space left, we might add a conference photograph here. Leave commented for now.
% \bookpartphoto[width=1.0\textwidth]{foobar.eps}{FooBar Photo (Photo: Any Photographer)}

\end{document}
